% Part: incompleteness
% Chapter: theories-computability
% Section: extensions-of-q-not-decidable

\documentclass[../../../include/open-logic-section]{subfiles}

\begin{document}

\olfileid[fr]{inc}{tcp}{cqn}

\olsection{Les extensions cohérentes de $\Th{Q}$ sont indécidables}

\begin{explain}
Rappelons qu'une théorie est \emph{cohérente} si elle ne
démontre, pour aucune !!{formula}~$!A$, à la fois
$!A$ et $\lnot !A$. Une contradiction permettant de
démontrer n'importe quel !!{sentence}, une théorie
incohérente est triviale : elle démontre tout
!!{sentence}. Toute théorie $\omega$-cohérente est
cohérente, mais la réciproque est fausse. Nous pouvons
donc affaiblir l'hypothèse du
\olref[oqn]{thm:oconsis-q} en ne supposant que la
cohérence, et obtenir un théorème plus fort.
\end{explain}

\begin{lem}
Il n'existe pas de « relation calculable universelle ».
Autrement dit, il n'existe aucune relation binaire
calculable $R(x,y)$ ayant la propriété suivante : pour
toute relation unaire calculable $S(y)$, il existe un
entier $k$ tel que, pour tout $y$, $S(y)$ est vraie si
et seulement si $R(k,y)$ est vraie.
\end{lem}

\begin{proof}
Supposons qu'une telle relation $R(x,y)$ existe.
La relation $S(y)$ définie par $\lnot R(y,y)$ est
calculable. Il existe donc un entier $k$ tel que
$S(y)$ équivaille à $R(k,y)$ pour tout $y$.
En particulier, $S(k)$ équivaut à la fois à
$R(k,k)$ et à $\lnot R(k,k)$, contradiction.
\end{proof}

\begin{thm}
Toute théorie cohérente $\Th{T}$ qui contient
$\Th{Q}$ est indécidable.
\end{thm}

\begin{proof}
Supposons que $\Th{T}$ soit une extension cohérente
et décidable de $\Th{Q}$. Nous allons en tirer une
relation calculable universelle, ce qui contredira
le lemme précédent.

Définissons $R(x,y)$ par la condition suivante :
\begin{quote}
$x$ code une !!{formula} $!D(u)$ sans variable libre autre que~$u$,
et $\Th{T}$ démontre $!D(\num y)$.
\end{quote}
La syntaxe et la condition sur les variables libres se vérifient
par calcul; la substitution de $\num y$ donne alors un !!{sentence}.
La décision supposée pour~$\Th{T}$ rend $R$ calculable; pour tout code invalide, $R$ est fausse. Montrons qu'elle
est universelle. Soit $S(y)$ une relation calculable
quelconque. Une !!{formula} $!D_S(u)$ la représente
dans $\Th{Q}$, donc aussi dans $\Th{T}$.
Pour chaque entier $n$, nous avons
\begin{eqnarray*}
S(n) & \liff & \Th{T}\vdash !D_S(\num n) \\
& \liff & R(\Gn{!D_S(u)},n),
\end{eqnarray*}
et
\begin{eqnarray*}
\lnot S(n) & \liff & \Th{T}\vdash\lnot !D_S(\num n) \\
& \liff & \Th{T}\not\vdash !D_S(\num n) \\
& \liff & \lnot R(\Gn{!D_S(u)},n).
\end{eqnarray*}
Les deux signes de la représentation, avec la cohérence
de~$\Th{T}$, donnent les premières équivalences.
Pour chaque~$n$, un des deux signes est démontrable
et l’autre ne l’est pas; les liens avec~$R$ suivent de sa définition.
Par conséquent,
pour tout $y$, $S(y)$ est vraie si et seulement si
$R(\Gn{!D_S(u)},y)$ est vraie. La relation $R$ est
universelle, contradiction.
\end{proof}

La « véritable arithmétique » est la théorie
$\Setabs{!A}{\Sat{\Nat}{!A}}$, c'est-à-dire
l'ensemble des !!{sentence}s arithmétiques vrais
dans l'interprétation standard.

\begin{cor}
La véritable arithmétique est indécidable.
\end{cor}

% La section sur l'indéfinissabilité de la vérité donnera un résultat
% plus fort, dû à Tarski.

\end{document}
